\documentclass{zhvt-classic}

% 使用本地源流明体（GenRyuMin2 TW）按文件路径加载，不依赖系统字体名
% 版框（内框／界栏／书口侧栏／鱼尾）改用与版心印章同一印泥棕色：
%   印章图片 yshs.png 实色 = RGB(135,68,52) = HTML 874434
\usepackage{xcolor}
\usepackage{tikz}
% sealbrown 双注册：xcolor 供 TikZ/pgf，L3 供 cls 版框
\definecolor{sealbrown}{HTML}{874434}
\ExplSyntaxOn
\color_set:nnn { sealbrown } { HTML } { 874434 }
\ExplSyntaxOff

\zhvtset{
  grid color       = sealbrown,
  main font        = GenRyuMin2TW-SB,
  main font path   = C:/Users/Administrator/texmf/fonts/opentype/public/GenRyuMin2TW/,
jiazhu font      = GenRyuMin2TW-M,
  jiazhu font path = C:/Users/Administrator/texmf/fonts/opentype/public/GenRyuMin2TW/,
  % 页码框容纳 4 字：正常页「九十七」+前置空 1 格=4 行恰好；
  % 100+ 页「一百零一」共 5 行超出，末字「一/二/三」被框裁掉（2026-10-04 实测），故放宽到 5
  shukou page number chars = 5,
  % 书口书名框高从默认 6 字减到 4 字，把省出的 2 字高拨给 chapter 框，
  % 使 6 字书名+卷名放进 9 字高的 chapter 框后不再压线（2026-10-05）
  shukou title chars = 4,
}

\title{史记}[秦本纪]
\author{汉·司马迁}
\maker{闲叟漫记·重製}


% ===== 索隐述赞样式（2026-11-08 四次定稿 · 愚大人亲手改定参数）=====
% ⚠ 竖排前提：自定义字体家族必须带 RawFeature={vertical:+vert}；
%   simhei.ttf 加 +vhal 会整体加载失败静默回退，只写 +vert；simfang 同理。
\newCJKfontfamily{\shuzanfont}[Path=C:/Windows/Fonts/,RawFeature={vertical:+vert}]{simhei.ttf}
\newCJKfontfamily{\shuzanbody}[Path=C:/Windows/Fonts/,RawFeature={vertical:+vert}]{simfang.ttf}
% 题头：152x30pt 圆角 6pt 棕底白字；baseline (0,2) 使整框落入界行格子不压版框线；
% 白字节点 x 不补偿、y +8.2pt（像素级墨迹实测 dx/dy≈0，L/R 留白约 20pt）。
\newcommand{\shuzanhead}[1]{%
  \begin{tikzpicture}[baseline={(0,2)}]
    \fill[sealbrown, rounded corners=6pt]
      (-76pt,-15pt) rectangle (76pt,15pt);
    \node[text=white, inner sep=0pt, anchor=center] at (0pt,8.2pt)
      {\shuzanfont\fontsize{23}{23}\selectfont #1};
  \end{tikzpicture}%
}
% ★ 修正书口卷名位置：chapter 框起始位置改回 cls 默认值，避免全局影响。
%   只在每个卷首页通过 \juan 命令里的空命令微调该页位置（安全，不会跨页污染）。
\ExplSyntaxOn
\AtBeginDocument{
  \dim_set:Nn \l__zhvt_shukou_chapter_title_start_dim
    { \l__zhvt_shukou_title_height_dim + \l__zhvt_font_size_dim * 2 }
  \dim_set:Nn \l__zhvt_shukou_page_number_start_dim
    { \l__zhvt_shukou_chapter_title_start_dim + \l__zhvt_shukou_chapter_title_height_dim }
}
\ExplSyntaxOff

\begin{document}

% ★★ 版心（书口）排布覆写：把书名放到鱼尾之下（2026-10-04）
%
%   cls 默认版心自上而下：书名（title 静态框，在鱼尾之上）→ 鱼尾 → 篇名（chapter 动态框）→ 页码
%   本模板改为：          鱼尾（位置不动）→ 书名 → 篇名 → 页码分割线 → 页码
%
%   手法：把 cls 放在 title 静态框里的书名，改放进 chapter 动态框（鱼尾之下）；
%         title 静态框留空。书名取 \l__zhvt_title_str，即 \title{...} 的第一参数，
%         所以换书时只改 \title，不必在此重复写书名。
%
%   ⚠ 三条必须遵守的写法（2026-10-04 实测）：
%     ① 汉字必须用 \hbox{...} 包住。flowfram 会把框内容装进一个宽 14.4mm（= 一个 line_sep
%        = 书口栏宽）的纵向 minipage 并按该宽度折行，而整页靠 PDF /Rotate 90 显示、
%        框内文字本身不变向 —— 不包 \hbox 时「围炉夜话」4 字会被逐字折成 4 行，
%        在版面上就成了横排 4 字并横向溢出版框。
%     ② 竖排下**一个 \hbox 就是一整列**。书名与篇名要同处一列，必须放进同一个 \hbox；
%        分成两个 \hbox 时，后一个会被当成新的一列、甩到隔壁栏去
%        （实测「序」跑到书口栏外侧 14.4mm 的正文列上）。中间的 \quad 在框内是横向空白，
%        显示出来正好是书名与篇名之间的纵向间隔。
%     ③ ★ 内容盒要写成 **与框等宽的定宽盒** `\hbox to \l__zhvt_shukou_chapter_title_height_dim{…\hfil}`。
%        实测：chapter 框是把内容盒沿纵向**居中**摆放的，内容宽度一变，书名就跟着漂：
%        自然宽度的 \hbox（书名 38.4mm）落在 x=157.2mm，加篇名后（57.6mm）落在 x=152.4mm，
%        目录页（无篇名）与正文页的书名就差半格。做成等宽盒后，无论有没有篇名，
%        内容固定从框顶起排，实测两页均落在 x=142.8mm，不再漂。
\ExplSyntaxOn
% 把书口 chapter 框（书名+篇名）起始位置下移 1 字高，避免 4 字卷名顶到鱼尾
% ⚠ cls 的 \__zhvt_calculate_dimensions: 在 \AtBeginDocument 里执行，会重算该 dim；
%   必须把覆写挂到 begin document 之后，否则会被二次覆盖（2026-10-05 实测）。
% (旧方案已移到导言区，直接修改 chapter 框起始位置)
% 卷标题命令 #1=书口篇名(鱼尾) #2=正文卷名（含卷序，卷名内空格用 ~）
%
% ★★ 目录页两项修订（2026-10-05，愚大人指出）★★
%   ① 页码起点统一对齐：
%      原状 —— 目录项写成 \par\hspace{4em}卷名\hspace{4em}页码，页码起点 = 卷名实宽 + 4em。
%      卷名实宽随字数变：「卷一~梁惠王上」6 字 +~、「卷七~离娄上」5 字 +~，
%      差 1 个字位 → 卷七~卷十 4 条的页码整体比其余 10 条高 27.21pt（一个字），看着上下不齐。
%      改法 —— 卷名补全角空格（宽 = 1 字 = 27.215pt）补足到 **7 字符位**再写进目录，
%      14 条页码起点遂全部落在同一条竖线上（起点对齐；末字随页码字数自然参差）。
%      ⚠ 只对目录里的副本 \l_tmpc_tl 补，正文卷名仍用 \l_tmpa_tl 原样 —— 正文像素不受影响。
%   ② 卷名整体上移半字位：前导 \hspace 由 4em 改 3.5em。
%      实测原状目录块在版框内偏下（天头留白 139.6pt / 地脚留白 110.7pt），
%      上移 0.5em = 13.6pt 后上下基本均分，视觉居中。
\NewDocumentCommand { \juan } { mm } {
	\clearpage
	\tl_set:Nn \l__zhvt_chapter_title_tl { #1 }
	\tl_set:Nn \l_tmpa_tl { #2 }
	\tl_set_eq:NN \l_tmpc_tl \l_tmpa_tl
	\int_step_inline:nn { 7 - \str_count:n { #2 } }
	  { \tl_put_right:Nn \l_tmpc_tl { 　 } }
	\tl_set:Nx \l_tmpb_tl { \zhnum { page } }
	\tl_put_left:Nn \l_tmpb_tl { \hspace { 4em } }
	\tl_put_left:Nn \l_tmpb_tl { \l_tmpc_tl }
	\tl_put_left:Nn \l_tmpb_tl { \par \hspace { 3.5em } }
	\addtocontents { toc } { \l_tmpb_tl }
	\skip_horizontal:n { \l__zhvt_font_size_dim * 4 }
	\l_tmpa_tl
	\par
	\thispagestyle { empty }
}
\cs_set_protected:Npn \__zhvt_set_shukou_content: {
	% ★ 鱼尾上方留白（title 静态框）嵌入闲叟钤印：
	%   整页 /Rotate 90，故此处把图用 \rotatebox{90} 转正为竖章；
	%   尺寸卡在版心条幅内（书内约 7mm 宽 × 33.8mm 高，留余量不压鱼尾）。
	\setstaticcontents * { title } {
		\vbox_to_ht:nn { \l__zhvt_baseline_skip_dim } {
			\vskip 14.2pt
			\hbox_to_wd:nn { \l__zhvt_shukou_title_height_dim } {
				\hfil
				\rotatebox [ origin = c ] { 90 } { \includegraphics [ width = 7mm , height = 33.8mm ] { yshs.png } }
				\hfil
			}
		}
	}
% ★ 版心（书口）统一改法（2026-11-16 愚大人定）：18pt＋框内居中＋书口文字同版框色
%   ① 原 27.2pt 时 7 字长篇名（如「廉颇蔺相如列传」）末字溢出 chapter 框、压住页码；
%      改 \fontsize{18pt}{18pt}，前后 \hfil 令字串在框内居中（不随篇名字数漂移）。
%   ② \vbox{\vskip 8.2pt} 定量下移，抵消缩字号后栏宽方向 8.2pt 上偏（像素实测 y 中心 543.3 = 原版）。
%      ⚠ \raisebox / \vspace* 在 flowfram 框内无效；换字号须同步换补偿：16pt→10.1、19pt→7.4。
%   ③ chapter 与 page~number 均置 \color{sealbrown}（HTML 874434），与版框／界行同色；正文仍黑。
%   ⚠ 页码字位、钤印、鱼尾、界行、正文一律未动；五部合订本同此口径。
	\setdynamiccontents * { chapter }
	{ \color{sealbrown} \vbox { \vskip 8.2pt \hbox to \l__zhvt_shukou_chapter_title_height_dim
		{ \hfil \fontsize{18pt}{18pt}\selectfont \l__zhvt_title_str \quad \l__zhvt_chapter_title_tl \hfil } } }
	\setdynamiccontents * { page~number } { \color{sealbrown} \quad \zhnum { page } }
}
\__zhvt_set_shukou_content:
\ExplSyntaxOff

\maketitle

% 单篇，无多卷目录，不渲染目录页

\juan{秦本纪}{卷五~秦本纪第五}
\hspace*{2em}秦之先。帝颛顼之苗裔孙曰女脩。女脩织。玄鸟陨卵。女脩吞之。生子大业。大业取少典之子。曰女华。女华生大费。与禹平水土。已成。帝锡玄圭。禹受曰。﹁非予能成。亦大费为辅。﹂帝舜曰。﹁咨尔费。赞禹功。其赐尔皂游。尔后嗣将大出。﹂乃妻之姚姓之玉女。大费拜受。佐舜调驯鸟兽。鸟兽多驯服。是为柏翳。舜赐姓嬴氏。\\
\hspace*{2em}大费生子二人。一曰大廉。实鸟俗氏。二曰若木。实费氏。其玄孙曰费昌。子孙或在中国。或在夷狄。费昌当夏桀之时。去夏归商。为汤御。以败桀于鸣条。大廉玄孙曰孟戏。中衍。鸟身人言。帝太戊闻而卜之使御。吉。遂致使御而妻之。自太戊以下。中衍之后。遂世有功。以佐殷国。故嬴姓多显。遂为诸侯。\\
\hspace*{2em}其玄孙曰中潏。在西戎。保西垂。生蜚廉。蜚廉生恶来。恶来有力。蜚廉善走。父子俱以材力事殷纣。周武王之伐纣。并杀恶来。是时蜚廉为纣石北方。还。无所报。为坛霍太山而报。得石棺。铭曰﹁帝令处父不与殷乱。赐尔石棺以华氏﹂。死。遂葬于霍太山。蜚廉复有子曰季胜。季胜生孟增。孟增幸于周成王。是为宅皋狼。皋狼生衡父。衡父生造父。造父以善御幸于周缪王。得骥。温骊。骅骝。騄耳之驷。西巡狩。乐而忘归。徐偃王作乱。造父为缪王御。长驱归周。一日千里以救乱。缪王以赵城封造父。造父族由此为赵氏。自蜚廉生季胜已下五世至造父。别居赵。赵衰其后也。恶来革者。蜚廉子也。蚤死。有子曰女防。女防生旁皋。旁皋生太几。太几生大骆。大骆生非子。以造父之宠。皆蒙赵城。姓赵氏。\\
\hspace*{2em}非子居犬丘。好马及畜。善养息之。犬丘人言之周孝王。孝王召使主马于汧渭之闲。马大蕃息。孝王欲以为大骆適嗣。申侯之女为大骆妻。生子成为適。申侯乃言孝王曰。﹁昔我先郦山之女。为戎胥轩妻。生中潏。以亲故归周。保西垂。西垂以其故和睦。今我复与大骆妻。生適子成。申骆重婚。西戎皆服。所以为王。王其图之。﹂于是孝王曰。﹁昔伯翳为舜主畜。畜多息。故有土。赐姓嬴。今其后世亦为朕息马。朕其分土为附庸。﹂邑之秦。使复续嬴氏祀。号曰秦嬴。亦不废申侯之女子为骆適者。以和西戎。\\
\hspace*{2em}秦嬴生秦侯。秦侯立十年。卒。生公伯。公伯立三年。卒。生秦仲。\\
\hspace*{2em}秦仲立三年。周厉王无道。诸侯或叛之。西戎反王室。灭犬丘大骆之族。周宣王即位。乃以秦仲为大夫。诛西戎。西戎杀秦仲。秦仲立二十三年。死于戎。有子五人。其长者曰庄公。周宣王乃召庄公昆弟五人。与兵七千人。使伐西戎。破之。于是复予秦仲后。及其先大骆地犬丘并有之。为西垂大夫。\\
\hspace*{2em}庄公居其故西犬丘。生子三人。其长男世父。世父曰。﹁戎杀我大父仲。我非杀戎王则不敢入邑。﹂遂将击戎。让其弟襄公。襄公为太子。庄公立四十四年。卒。太子襄公代立。襄公元年。以女弟缪嬴为丰王妻。襄公二年。戎围犬丘。世父击之。为戎人所虏。岁余。复归世父。七年春。周幽王用褒姒废太子。立褒姒子为嫡。数欺诸侯。诸侯叛之。西戎犬戎与申侯伐周。杀幽王郦山下。而秦襄公将兵救周。战甚力。有功。周避犬戎难。东徙雒邑。襄公以兵送周平王。平王封襄公为诸侯。赐之岐以西之地。曰。﹁戎无道。侵夺我岐。丰之地。秦能攻逐戎。即有其地。﹂与誓。封爵之。襄公于是始国。与诸侯通使聘享之礼。乃用骝驹。黄牛。羝羊各三。祠上帝西畤。十二年。伐戎而至岐。卒。生文公。\\
\hspace*{2em}文公元年。居西垂宫。三年。文公以兵七百人东猎。四年。至汧渭之会。曰。﹁昔周邑我先秦嬴于此。后卒获为诸侯。﹂乃卜居之。占曰吉。即营邑之。十年。初为鄜畤。用三牢。十三年。初有史以纪事。民多化者。十六年。文公以兵伐戎。戎败走。于是文公遂收周余民有之。地至岐。岐以东献之周。十九年。得陈宝。二十年。法初有三族之罪。二十七年。伐南山大梓。丰大特。四十八年。文公太子卒。赐谥为竫公。竫公之长子为太子。是文公孙也。五十年。文公卒。葬西山。竫公子立。是为宁公。\\
\hspace*{2em}宁公二年。公徙居平阳。遣兵伐荡社。三年。与亳战。亳王奔戎。遂灭荡社。四年。鲁公子翚弑其君隐公。十二年。伐荡氏。取之。宁公生十岁立。立十二年卒。葬西山。生子三人。长男武公为太子。武公弟德公。同母鲁姬子。生出子。宁公卒。大庶长弗忌。威垒。三父废太子而立出子为君。出子六年。三父等复共令人贼杀出子。出子生五岁立。立六年卒。三父等乃复立故太子武公。\\
\hspace*{2em}武公元年。伐彭戏氏。至于华山下。居平阳封宫。三年。诛三父等而夷三族。以其杀出子也。郑高渠眯杀其君昭公。十年。伐邽。冀戎。初县之。十一年。初县杜。郑。灭小虢。\\
\hspace*{2em}十三年。齐人管至父。连称等杀其君襄公而立公孙无知。晋灭霍。魏。耿。齐雍廪杀无知。管至父等而立齐桓公。齐。晋为强国。\\
\hspace*{2em}十九年。晋曲沃始为晋侯。齐桓公伯于鄄。\\
\hspace*{2em}二十年。武公卒。葬雍平阳。初以人从死。从死者六十六人。有子一人。名曰白。白不立。封平阳。立其弟德公。\\
\hspace*{2em}德公元年。初居雍城大郑宫。以牺三百牢祠鄜畤。卜居雍。后子孙饮马于河。梁伯。芮伯来朝。二年。初伏。以狗御蛊。德公生三十三岁而立。立二年卒。生子三人。长子宣公。中子成公。少子穆公。长子宣公立。\\
\hspace*{2em}宣公元年。卫。燕伐周。出惠王。立王子穨。三年。郑伯。虢叔杀子穨而入惠王。四年。作密畤。与晋战河阳。胜之。十二年。宣公卒。生子九人。莫立。立其弟成公。\\
\hspace*{2em}成公元年。梁伯。芮伯来朝。齐桓公伐山戎。次于孤竹。\\
\hspace*{2em}成公立四年卒。子七人。莫立。立其弟缪公。\\
\hspace*{2em}缪公任好元年。自将伐茅津。胜之。四年。迎妇于晋。晋太子申生姊也。其岁。齐桓公伐楚。至邵陵。\\
\hspace*{2em}五年。晋献公灭虞。虢。虏虞君与其大夫百里傒。以璧马赂于虞故也。既虏百里傒。以为秦缪公夫人媵于秦。百里傒亡秦走宛。楚鄙人执之。缪公闻百里傒贤。欲重赎之。恐楚人不与。乃使人谓楚曰。﹁吾媵臣百里傒在焉。请以五羖羊皮赎之。﹂。楚人遂许与之。当是时。百里傒年已七十余。缪公释其囚。与语国事。谢曰。﹁臣亡国之臣。何足问。﹂缪公曰。﹁虞君不用子。故亡。非子罪也。﹂固问。语三日。缪公大说。授之国政。号曰五羖大夫。百里傒让曰。﹁臣不及臣友蹇叔。蹇叔贤而世莫知。臣常游困于齐而乞食铚人。蹇叔收臣。臣因而欲事齐君无知。蹇叔止臣。臣得脱齐难。遂之周。周王子穨好牛。臣以养牛干之。及穨欲用臣。蹇叔止臣。臣去。得不诛。事虞君。蹇叔止臣。臣知虞君不用臣。臣诚私利禄爵。且留。再用其言。得脱。一不用。及虞君难。是以知其贤。﹂于是缪公使人厚币迎蹇叔。以为上大夫。\\
\hspace*{2em}秋。缪公自将伐晋。战于河曲。晋骊姬作乱。太子申生死新城。重耳。夷吾出奔。\\
\hspace*{2em}九年。齐桓公会诸侯于葵丘。\\
\hspace*{2em}晋献公卒。立骊姬子奚齐。其臣里克杀奚齐。荀息立卓子。克又杀卓子及荀息。夷吾使人请秦。求入晋。于是缪公许之。使百里傒将兵送夷吾。夷吾谓曰。﹁诚得立。请割晋之河西八城与秦。﹂及至。已立。而使丕郑谢秦。背约不与河西城。而杀里克。丕郑闻之。恐。因与缪公谋曰。﹁晋人不欲夷吾。实欲重耳。今背秦约而杀里克。皆吕甥。郄芮之计也。愿君以利急召吕。郄。吕。郄至。则更入重耳便。﹂缪公许之。使人与丕郑归。召吕。郄。吕。郄等疑丕郑有间。乃言夷吾杀丕郑。丕郑子丕豹奔秦。说缪公曰。﹁晋君无道。百姓不亲。可伐也。﹂缪公曰。﹁百姓苟不便。何故能诛其大臣。能诛其大臣。此其调也。﹂不听。而阴用豹。\\
\hspace*{2em}十二年。齐管仲。隰朋死。\\
\hspace*{2em}晋旱。来请粟。丕豹说缪公勿与。因其饥而伐之。缪公问公孙支。支曰。﹁饥穰更事耳。不可不与。﹂问百里傒。傒曰。﹁夷吾得罪于君。其百姓何罪。﹂于是用百里傒。公孙支言。卒与之粟。以船漕车转。自雍相望至绛。\\
\hspace*{2em}十四年。秦饥。请粟于晋。晋君谋之群臣。虢射曰。﹁因其饥伐之。可有大功。﹂晋君从之。十五年。兴兵将攻秦。缪公发兵。使丕豹将。自往击之。九月壬戌。与晋惠公夷吾合战于韩地。晋君弃其军。与秦争利。还而马騺。缪公与麾下驰追之。不能得晋君。反为晋军所围。晋击缪公。缪公伤。于是岐下食善马者三百人驰冒晋军。晋军解围。遂脱缪公而反生得晋君。初。缪公亡善马。岐下野人共得而食之者三百余人。吏逐得。欲法之。缪公曰。﹁君子不以畜产害人。吾闻食善马肉不饮酒。伤人。﹂乃皆赐酒而赦之。三百人者闻秦击晋。皆求从。从而见缪公窘。亦皆推锋争死。以报食马之德。于是缪公虏晋君以归。令于国。齐宿。吾将以晋君祠上帝。周天子闻之。曰﹁晋我同姓﹂。为请晋君。夷吾姊亦为缪公夫人。夫人闻之。乃衰绖跣。曰。﹁妾兄弟不能相救。以辱君命。﹂缪公曰。﹁我得晋君以为功。今天子为请。夫人是忧。﹂乃与晋君盟。许归之。更舍上舍。而馈之七牢。十一月。归晋君夷吾。夷吾献其河西地。使太子圉为质于秦。秦妻子圉以宗女。是时秦地东至河。\\
\hspace*{2em}十八年。齐桓公卒。二十年。秦灭梁。芮。\\
\hspace*{2em}二十二年。晋公子圉闻晋君病。曰。﹁梁。我母家也。而秦灭之。我兄弟多。即君百岁后。秦必留我。而晋轻。亦更立他子。﹂子圉乃亡归晋。二十三年。晋惠公卒。子圉立为君。秦怨圉亡去。乃迎晋公子重耳于楚。而妻以故子圉妻。重耳初谢。后乃受。缪公益礼厚遇之。二十四年春。秦使人告晋大臣。欲入重耳。晋许之。于是使人送重耳。二月。重耳立为晋君。是为文公。文公使人杀子圉。子圉是为怀公。\\
\hspace*{2em}其秋。周襄王弟带以翟伐王。王出居郑。二十五年。周王使人告难于晋。秦。秦缪公将兵助晋文公入襄王。杀王弟带。二十八年。晋文公败楚于城濮。三十年。缪公助晋文公围郑。郑使人言缪公曰。﹁亡郑厚晋。于晋而得矣。而秦未有利。晋之强。秦之忧也。﹂缪公乃罢兵归。晋亦罢。三十二年冬。晋文公卒。\\
\hspace*{2em}郑人有卖郑于秦曰。﹁我主其城门。郑可袭也。﹂缪公问蹇叔。百里傒。对曰。﹁径数国千里而袭人。希有得利者。且人卖郑。庸知我国人不有以我情告郑者乎。不可。﹂缪公曰。﹁子不知也。吾已决矣。﹂遂发兵。使百里傒子孟明视。蹇叔子西乞术及白乙丙将兵。行日。百里傒。蹇叔二人哭之。缪公闻。怒曰。﹁孤发兵而子沮哭吾军。何也。﹂二老曰。﹁臣非敢沮君军。军行。臣子与往。臣老。迟还恐不相见。故哭耳。﹂二老退。谓其子曰。﹁汝军即败。必于肴阨矣。﹂三十三年春。秦兵遂东。更晋地。过周北门。周王孙满曰。﹁秦师无礼。不败何待。﹂兵至滑。郑贩卖贾人弦高。持十二牛将卖之周。见秦兵。恐死虏。因献其牛。曰。﹁闻大国将诛郑。郑君谨修守御备。使臣以牛十二劳军士。﹂秦三将军相谓曰。﹁将袭郑。郑今已觉之。往无及已。﹂灭滑。滑。晋之边邑也。\\
\hspace*{2em}当是时。晋文公丧尚未葬。太子襄公怒曰。﹁秦侮我孤。因丧破我滑。﹂遂墨衰绖。发兵遮秦兵于肴。击之。大破秦军。无一人得脱者。虏秦三将以归。文公夫人。秦女也。为秦三囚将请曰。﹁缪公之怨此三人入于骨髓。愿令此三人归。令我君得自快烹之。﹂晋君许之。归秦三将。三将至。缪公素服郊迎。向三人哭曰。﹁孤以不用百里傒。蹇叔言以辱三子。三子何罪乎。子其悉心雪耻。毋怠。﹂遂复三人官秩如故。愈益厚之。\\
\hspace*{2em}三十四年。楚太子商臣弑其父成王代立。\\
\hspace*{2em}缪公于是复使孟明视等将兵伐晋。战于彭衙。秦不利。引兵归。\\
\hspace*{2em}戎王使由余于秦。由余。其先晋人也。亡入戎。能晋言。闻缪公贤。故使由余观秦。秦缪公示以宫室。积聚。由余曰。﹁使鬼为之。则劳神矣。使人为之。亦苦民矣。﹂缪公怪之。问曰。﹁中国以诗书礼乐法度为政。然尚时乱。今戎夷无此。何以为治。不亦难乎。﹂由余笑曰。﹁此乃中国所以乱也。夫自上圣黄帝作为礼乐法度。身以先之。仅以小治。及其后世。日以骄淫。阻法度之威。以责督于下。下罢极则以仁义怨望于上。上下交争怨而相篡弑。至于灭宗。皆以此类也。夫戎夷不然。上含淳德以遇其下。下怀忠信以事其上。一国之政犹一身之治。不知所以治。此真圣人之治也。﹂于是缪公退而问内史廖曰。﹁孤闻邻国有圣人。敌国之忧也。今由余贤。寡人之害。将奈之何。﹂内史廖曰。﹁戎王处辟匿。未闻中国之声。君试遗其女乐。以夺其志。为由余请。以疏其间。留而莫遣。以失其期。戎王怪之。必疑由余。君臣有间。乃可虏也。且戎王好乐。必怠于政。﹂缪公曰。﹁善。﹂因与由余曲席而坐。传器而食。问其地形与其兵势尽察。而后令内史廖以女乐二八遗戎王。戎王受而说之。终年不还。于是秦乃归由余。由余数谏不听。缪公又数使人间要由余。由余遂去降秦。缪公以客礼礼之。问伐戎之形。\\
\hspace*{2em}三十六年。缪公复益厚孟明等。使将兵伐晋。渡河焚船。大败晋人。取王官及鄗。以报肴之役。晋人皆城守不敢出。于是缪公乃自茅津渡河。封肴中尸。为发丧。哭之三日。乃誓于军曰。﹁嗟士卒。听无哗。余誓告汝。古之人谋黄髪番番。则无所过。﹂以申思不用蹇叔。百里傒之谋。故作此誓。令后世以记余过。君子闻之。皆为垂涕。曰。﹁嗟乎。秦缪公之与人周也。卒得孟明之庆。﹂\\
\hspace*{2em}三十七年。秦用由余谋伐戎王。益国十二。开地千里。遂霸西戎。天子使召公过贺缪公以金鼓。三十九年。缪公卒。葬雍。从死者百七十七人。秦之良臣子舆氏三人名曰奄息。仲行。针虎。亦在从死之中。秦人哀之。为作歌黄鸟之诗。君子曰。﹁秦缪公广地益国。东服强晋。西霸戎夷。然不为诸侯盟主。亦宜哉。死而弃民。收其良臣而从死。且先王崩。尚犹遗德垂法。况夺之善人良臣百姓所哀者乎。是以知秦不能复东征也。﹂缪公子四十人。其太子嵤代立。是为康公。\\
\hspace*{2em}康公元年。往岁缪公之卒。晋襄公亦卒。襄公之弟名雍。秦出也。在秦。晋赵盾欲立之。使随会来迎雍。秦以兵送至令狐。晋立襄公子而反击秦师。秦师败。随会来奔。二年。秦伐晋。取武城。报令狐之役。四年。晋伐秦。取少梁。六年。秦伐晋。取羁马。战于河曲。大败晋军。晋人患随会在秦为乱。乃使魏雠馀详反。合谋会。诈而得会。会遂归晋。康公立十二年卒。子共公立。\\
\hspace*{2em}共公二年。晋赵穿弑其君灵公。三年。楚庄王强。北兵至雒。问周鼎。共公立五年卒。子桓公立。\\
\hspace*{2em}桓公三年。晋败我一将。十年。楚庄王服郑。北败晋兵于河上。当是之时。楚霸。为会盟合诸侯。二十四年。晋厉公初立。与秦桓公夹河而盟。归而秦倍盟。与翟合谋击晋。二十六年。晋率诸侯伐秦。秦军败走。追至泾而还。桓公立二十七年卒。子景公立。\\
\hspace*{2em}景公四年。晋栾书弑其君厉公。十五年。救郑。败晋兵于栎。是时晋悼公为盟主。十八年。晋悼公强。数会诸侯。率以伐秦。败秦军。秦军走。晋兵追之。遂渡泾。至棫林而还。二十七年。景公如晋。与平公盟。已而背之。三十六年。楚公子围弑其君而自立。是为灵王。景公母弟后子针有宠。景公母弟富。或谮之。恐诛。乃奔晋。车重千乘。晋平公曰。﹁后子富如此。何以自亡。﹂对曰。﹁秦公无道。畏诛。欲待其后世乃归。﹂三十九年。楚灵王强。会诸侯于申。为盟主。杀齐庆封。景公立四十年卒。子哀公立。后子复来归秦。\\
\hspace*{2em}哀公八年。楚公子弃疾弑灵王而自立。是为平王。十一年。楚平王来求秦女为太子建妻。至国。女好而自娶之。十五年。楚平王欲诛建。建亡。伍子胥奔吴。晋公室卑而六卿强。欲内相攻。是以久秦晋不相攻。三十一年。吴王阖闾与伍子胥伐楚。楚王亡奔随。吴遂入郢。楚大夫申包胥来告急。七日不食。日夜哭泣。于是秦乃发五百乘救楚。败吴师。吴师归。楚昭王乃得复入郢。哀公立三十六年卒。太子夷公。夷公蚤死。不得立。立夷公子。是为惠公。\\
\hspace*{2em}惠公元年。孔子行鲁相事。五年。晋卿中行。范氏反晋。晋使智氏。赵简子攻之。范。中行氏亡奔齐。惠公立十年卒。子悼公立。\\
\hspace*{2em}悼公二年。齐臣田乞弑其君孺子。立其兄阳生。是为悼公。六年。吴败齐师。齐人弑悼公。立其子简公。九年。晋定公与吴王夫差盟。争长于黄池。卒先吴。吴强。陵中国。十二年。齐田常弑简公。立其弟平公。常相之。十三年。楚灭陈。秦悼公立十四年卒。子厉共公立。孔子以悼公十二年卒。\\
\hspace*{2em}厉共公二年。蜀人来赂。十六年。堑河旁。以兵二万伐大荔。取其王城。二十一年。初县频阳。晋取武成。二十四年。晋乱。杀智伯。分其国与赵。韩。魏。二十五年。智开与邑人来奔。三十三年。伐义渠。虏其王。三十四年。日食。厉共公卒。子躁公立。躁公二年。南郑反。十三年。义渠来伐。至渭南。十四年。躁公卒。立其弟怀公。\\
\hspace*{2em}怀公四年。庶长晁与大臣围怀公。怀公自杀。怀公太子曰昭子。蚤死。大臣乃立太子昭子之子。是为灵公。灵公。怀公孙也。\\
\hspace*{2em}灵公六年。晋城少梁。秦击之。十三年。城籍姑。灵公卒。子献公不得立。立灵公季父悼子。是为简公。简公。昭子之弟而怀公子也。\\
\hspace*{2em}简公六年。令吏初带剑。堑洛。城重泉。十六年卒。子惠公立。\\
\hspace*{2em}惠公十二年。子出子生。十三年。伐蜀。取南郑。惠公卒。出子立。\\
\hspace*{2em}出子二年。庶长改迎灵公之子献公于河西而立之。杀出子及其母。沈之渊旁。秦以往者数易君。君臣乖乱。故晋复强。夺秦河西地。\\
\hspace*{2em}献公元年。止从死。二年。城栎阳。四年正月庚寅。孝公生。十一年。周太史儋见献公曰。﹁周故与秦国合而别。别五百岁复合。合十七岁而霸王出。﹂十六年。桃冬花。十八年。雨金栎阳。二十一年。与晋战于石门。斩首六万。天子贺以黼黻。二十三年。与魏晋战少梁。虏其将公孙痤。二十四年。献公卒。子孝公立。年已二十一岁矣。\\
\hspace*{2em}孝公元年。河山以东强国六。与齐威。楚宣。魏惠。燕悼。韩哀。赵成侯并。淮泗之间小国十余。楚。魏与秦接界。魏筑长城。自郑滨洛以北。有上郡。楚自汉中。南有巴。黔中。周室微。诸侯力政。争相并。秦僻在雍州。不与中国诸侯之会盟。夷翟遇之。孝公于是布惠。振孤寡。招战士。明功赏。下令国中曰。﹁昔我缪公自岐雍之间。修德行武。东平晋乱。以河为界。西霸戎翟。广地千里。天子致伯。诸侯毕贺。为后世开业。甚光美。会往者厉。躁。简公。出子之不宁。国家内忧。未遑外事。三晋攻夺我先君河西地。诸侯卑秦。丑莫大焉。献公即位。镇抚边境。徙治栎阳。且欲东伐。复缪公之故地。修缪公之政令。寡人思念先君之意。常痛于心。宾客群臣有能出奇计强秦者。吾且尊官。与之分土。﹂于是乃出兵东围陕城。西斩戎之獂王。\\
\hspace*{2em}卫鞅闻是令下。西入秦。因景监求见孝公。\\
\hspace*{2em}二年。天子致胙。\\
\hspace*{2em}三年。卫鞅说孝公变法修刑。内务耕稼。外劝战死之赏罚。孝公善之。甘龙。杜挚等弗然。相与争之。卒用鞅法。百姓苦之。居三年。百姓便之。乃拜鞅为左庶长。其事在商君语中。\\
\hspace*{2em}七年。与魏惠王会杜平。八年。与魏战元里。有功。十年。卫鞅为大良造。将兵围魏安邑。降之。十二年。作为咸阳。筑冀阙。秦徙都之。并诸小乡聚。集为大县。县一令。四十一县。为田开阡陌。东地渡洛。十四年。初为赋。十九年。天子致伯。二十年。诸侯毕贺。秦使公子少官率师会诸侯逢泽。朝天子。\\
\hspace*{2em}二十一年。齐败魏马陵。\\
\hspace*{2em}二十二年。卫鞅击魏。虏魏公子昂。封鞅为列侯。号商君。\\
\hspace*{2em}二十四年。与晋战雁门。虏其将魏错。\\
\hspace*{2em}孝公卒。子惠文君立。是岁。诛卫鞅。鞅之初为秦施法。法不行。太子犯禁。鞅曰。﹁法之不行。自于贵戚。君必欲行法。先于太子。太子不可黥。黥其傅师。﹂于是法大用。秦人治。及孝公卒。太子立。宗室多怨鞅。鞅亡。因以为反。而卒车裂以徇秦国。\\
\hspace*{2em}惠文君元年。楚。韩。赵。蜀人来朝。二年。天子贺。三年。王冠。四年。天子致文武胙。齐。魏为王。\\
\hspace*{2em}五年。阴晋人犀首为大良造。六年。魏纳阴晋。阴晋更名宁秦。七年。公子昂与魏战。虏其将龙贾。斩首八万。八年。魏纳河西地。九年。渡河。取汾阴。皮氏。与魏王会应。围焦。降之。十年。张仪相秦。魏纳上郡十五县。十一年。县义渠。归魏焦。曲沃。义渠君为臣。更名少梁曰夏阳。十二年。初腊。十三年四月戊午。魏君为王。韩亦为王。使张仪伐取陕。出其人与魏。\\
\hspace*{2em}十四年。更为元年。二年。张仪与齐。楚大臣会啮桑。三年。韩。魏太子来朝。张仪相魏。五年。王游至北河。七年。乐池相秦。韩。赵。魏。燕。齐帅匈奴共攻秦。秦使庶长疾与战修鱼。虏其将申差。败赵公子渴。韩太子奂。斩首八万二千。八年。张仪复相秦。九年。司马错伐蜀。灭之。伐取赵中都。西阳。十年。韩太子苍来质。伐取韩石章。伐败赵将泥。伐取义渠二十五城。十一年。樗里疾攻魏焦。降之。败韩岸门。斩首万。其将犀首走。公子通封于蜀。燕君让其臣子之。十二年。王与梁王会临晋。庶长疾攻赵。虏赵将庄。张仪相楚。十三年。庶长章击楚于丹阳。虏其将屈丐。斩首八万。又攻楚汉中。取地六百里。置汉中郡。楚围雍氏。秦使庶长疾助韩而东攻齐。到满助魏攻燕。十四年。伐楚。取召陵。丹。犁臣。蜀相壮杀蜀侯来降。\\
\hspace*{2em}惠王卒。子武王立。韩。魏。齐。楚。越皆宾从。\\
\hspace*{2em}武王元年。与魏襄王会临晋。诛蜀相壮。张仪。魏章皆东出之魏。伐义渠。丹。犁。二年。初置丞相。樗里疾。甘茂为左右丞相。张仪死于魏。三年。与韩襄王会临晋外。南公揭卒。樗里疾相韩。武王谓甘茂曰。﹁寡人欲容车通三川。窥周室。死不恨矣。﹂其秋。使甘茂。庶长封伐宜阳。四年。拔宜阳。斩首六万。涉河。城武遂。魏太子来朝。武王有力好戏。力士任鄙。乌获。孟说皆至大官。王与孟说举鼎。绝膑。八月。武王死。族孟说。武王取魏女为后。无子。立异母弟。是为昭襄王。昭襄母楚人。姓芈氏。号宣太后。武王死时。昭襄王为质于燕。燕人送归。得立。\\
\hspace*{2em}昭襄王元年。严君疾为相。甘茂出之魏。二年。彗星见。庶长壮与大臣。诸侯。公子为逆。皆诛。及惠文后皆不得良死。悼武王后出归魏。三年。王冠。与楚王会黄棘。与楚上庸。四年。取蒲阪。彗星见。五年。魏王来朝应亭。复与魏蒲阪。六年。蜀侯煇反。司马错定蜀。庶长奂伐楚。斩首二万。泾阳君质于齐。日食。昼晦。七年。拔新城。樗里子卒。八年。使将军芈戎攻楚。取新市。齐使章子。魏使公孙喜。韩使暴鸢共攻楚方城。取唐眛。赵破中山。其君亡。竟死齐。魏公子劲。韩公子长为诸侯。九年。孟尝君薛文来相秦。奂攻楚。取八城。杀其将景快。十年。楚怀王入朝秦。秦留之。薛文以金受免。楼缓为丞相。十一年。齐。韩。魏。赵。宋。中山五国共攻秦。至盐氏而还。秦与韩。魏河北及封陵以和。彗星见。楚怀王走之赵。赵不受。还之秦。即死。归葬。十二年。楼缓免。穰侯魏冉为相。予楚粟五万石。\\
\hspace*{2em}十三年。向寿伐韩。取武始。左更白起攻新城。五大夫礼出亡奔魏。任鄙为汉中守。十四年。左更白起攻韩。魏于伊阙。斩首二十四万。虏公孙喜。拔五城。十五年。大良造白起攻魏。取垣。复予之。攻楚。取宛。十六年。左更错取轵及邓。冉免。封公子市宛。公子悝邓。魏冉陶。为诸侯。十七年。城阳君入朝。及东周君来朝。秦以垣为蒲阪。皮氏。王之宜阳。十八年。错攻垣。河雍。决桥取之。十九年。王为西帝。齐为东帝。皆复去之。吕礼来自归。齐破宋。宋王在魏。死温。任鄙卒。二十年。王之汉中。又之上郡。北河。二十一年。错攻魏河内。魏献安邑。秦出其人。募徙河东赐爵。赦罪人迁之。泾阳君封宛。二十二年。蒙武伐齐。河东为九县。与楚王会宛。与赵王会中阳。二十三年。尉斯离与三晋。燕伐齐。破之济西。王与魏王会宜阳。与韩王会新城。二十四年。与楚王会鄢。又会穰。秦取魏安城。至大梁。燕。赵救之。秦军去。魏冉免相。二十五年。拔赵二城。与韩王会新城。与魏王会新明邑。二十六年。赦罪人迁之穰。侯冉复相。二十七年。错攻楚。赦罪人迁之南阳。白起攻赵。取代光狼城。又使司马错发陇西。因蜀攻楚黔中。拔之。二十八年。大良造白起攻楚。取鄢。邓。赦罪人迁之。二十九年。大良造白起攻楚。取郢为南郡。楚王走。周君来。王与楚王会襄陵。白起为武安君。三十年。蜀守若伐楚。取巫郡。及江南为黔中郡。三十一年。白起伐魏。取两城。楚人反我江南。三十二年。相穰侯攻魏。至大梁。破暴鸢。斩首四万。鸢走。魏入三县请和。三十三年。客卿胡阳攻魏卷。蔡阳。长社。取之。击芒卯华阳。破之。斩首十五万。魏入南阳以和。三十四年。秦与魏。韩上庸地为一郡。南阳免臣迁居之。三十五年。佐韩。魏。楚伐燕。初置南阳郡。三十六年。客卿灶攻齐。取刚。寿。予穰侯。三十八年。中更胡阳攻赵阏与。不能取。四十年。悼太子死魏。归葬芷阳。四十一年夏。攻魏。取邢丘。怀。四十二年。安国君为太子。十月。宣太后薨。葬芷阳郦山。九月。穰侯出之陶。四十三年。武安君白起攻韩。拔九城。斩首五万。四十四年。攻韩南阳。取之。四十五年。五大夫贲攻韩。取十城。叶阳君悝出之国。未至而死。四十七年。秦攻韩上党。上党降赵。秦因攻赵。赵发兵击秦。相距。秦使武安君白起击。大破赵于长平。四十余万尽杀之。四十八年十月。韩献垣雍。秦军分为三军。武安君归。王龁将伐赵武安。皮牢。拔之。司马梗北定太原。尽有韩上党。正月。兵罢。复守上党。其十月。五大夫陵攻赵邯郸。四十九年正月。益发卒佐陵。陵战不善。免。王龁代将。其十月。将军张唐攻魏。为蔡尉捐弗守。还斩之。五十年十月。武安君白起有罪。为士伍。迁阴密。张唐攻郑。拔之。十二月。益发卒军汾城旁。武安君白起有罪。死。龁攻邯郸。不拔。去。还奔汾军二月余。攻晋军。斩首六千。晋楚流死河二万人。攻汾城。即从唐拔宁新中。宁新中更名安阳。初作河桥。\\
\hspace*{2em}五十一年。将军摎攻韩。取阳城。负黍。斩首四万。攻赵。取二十余县。首虏九万。西周君背秦。与诸侯约从。将天下锐兵出伊阙攻秦。令秦毋得通阳城。于是秦使将军摎攻西周。西周君走来自归。顿首受罪。尽献其邑三十六城。口三万。秦王受献。归其君于周。五十二年。周民东亡。其器九鼎入秦。周初亡。\\
\hspace*{2em}五十三年。天下来宾。魏后。秦使摎伐魏。取吴城。韩王入朝。魏委国听令。五十四年。王郊见上帝于雍。五十六年秋。昭襄王卒。子孝文王立。尊唐八子为唐太后。而合其葬于先王。韩王衰绖入吊祠。诸侯皆使其将相来吊祠。视丧事。\\
\hspace*{2em}孝文王元年。赦罪人。修先王功臣。褒厚亲戚。弛苑囿。孝文王除丧。十月己亥即位。三日辛丑卒。子庄襄王立。\\
\hspace*{2em}庄襄王元年。大赦罪人。修先王功臣。施德厚骨肉而布惠于民。东周君与诸侯谋秦。秦使相国吕不韦诛之。尽入其国。秦不绝其祀。以阳人地赐周君。奉其祭祀。使蒙骜伐韩。韩献成皋。巩。秦界至大梁。初置三川郡。二年。使蒙骜攻赵。定太原。三年。蒙骜攻魏高都。汲。拔之。攻赵榆次。新城。狼孟。取三十七城。四月日食。王龁攻上党。初置太原郡。魏将无忌率五国兵击秦。秦却于河外。蒙骜败。解而去。五月丙午。庄襄王卒。子政立。是为秦始皇帝。\\
\hspace*{2em}秦王政立二十六年。初并天下为三十六郡。号为始皇帝。始皇帝五十一年而崩。子胡亥立。是为二世皇帝。三年。诸侯并起叛秦。赵高杀二世。立子婴。子婴立月余。诸侯诛之。遂灭秦。其语在始皇本纪中。\\
\hspace*{2em}太史公曰。秦之先为嬴姓。其后分封。以国为姓。有徐氏。郯氏。莒氏。终黎氏。运奄氏。菟裘氏。将梁氏。黄氏。江氏。修鱼氏。白冥氏。蜚廉氏。秦氏。然秦以其先造父封赵城。为赵氏。\\

% ===== 索隐述赞（司马贞《史记索隐》）· 统一样式 =====
% 题头「索隐述赞曰」五字并入反白框；正文每四句一列，定宽盒 \hfil 居列中；
% \fontsize{24}{41} 列距对齐全书界行（121+41k）；\clearpage 使述赞独占一页。
\clearpage
\shuzanhead{索隐述赞曰}\linebreak
{\shuzanbody\fontsize{24}{41}\selectfont
\hbox to \hsize{\hfil 柏翳佐舜。皂斿是旌。蜚廉事纣。石椁斯营。\hfil}\linebreak
\hbox to \hsize{\hfil 造父善驭。封之赵城。非子息马。厥号秦嬴。\hfil}\linebreak
\hbox to \hsize{\hfil 礼乐射御。西垂有声。襄公救周。始命列国。\hfil}\linebreak
\hbox to \hsize{\hfil 金祠白帝。龙祚水德。祥应陈宝。妖除丰特。\hfil}\linebreak
\hbox to \hsize{\hfil 里奚致霸。卫鞅任刻。厥后吞并。卒成凶慝。\hfil}%
\par}
\end{document}
